\subsection{Pseudospectrum}

DEFINITION 9. --- Let E be a complex Banach space and let u be a closed partial operator on E. Let \(\varepsilon\) be a strictly positive real number. We call the \(\varepsilon\)-pseudospectrum of u the union of the spectrum of u and the set of complex numbers \(\lambda\) belonging to the resolvent set of u such that \(\| \mathrm{R}(u,\lambda)\| >\varepsilon^{-1}\) . We denote this set by \(\mathrm{PSp}_{\varepsilon}(u)\) .

\pandocbounded{\includegraphics[keepaspectratio]{images/part002_14758b8b.jpg}}

Some authors define the \(\varepsilon\) -pseudo-spectrum of \(u\) as the set \(\mathrm{T}_{\varepsilon}(u)\), which is the union of \(\mathrm{Sp}(u)\) and the set of \(\lambda \in \mathbf{C} - \mathrm{Sp}(u)\) such that \(\| \mathrm{R}(u,\lambda)\| \geqslant \varepsilon^{-1}\) . The closure of \(\mathrm{PSp}_{\varepsilon}(u)\) is contained in \(\mathrm{T}_{\varepsilon}(u)\) but this inclusion can be strict, even if E is a Hilbert space (cf. exercises 18 of IV, p. 348 and 19 of IV, p. 349).

PROPOSITION 18. --- Let E be a complex Banach space and u a closed partial operator on E. Let \(\varepsilon \in \mathbf{R}_{+}^{*}\) . The set \(\mathrm{PSp}_{\varepsilon}(u)\) is an open subset of C. It is nonempty if E is nonzero.

If E is zero, then \(\mathrm{PSp}_{\varepsilon}(u)\) is empty. Suppose that E is nonzero. By prop. 14 of IV, p. 246, the set \(\mathrm{PSp}_{\varepsilon}(u)\) is an open subset of C.

The set \(\mathrm{PSp}_{\varepsilon}(u)\) is not empty if the spectrum of \(u\) is not empty. If \(\mathrm{Sp}(u)\) is empty, then Liouville's theorem (VAR, R1, p. 29, 3.3.6) implies that the holomorphic function on \(\mathbf{C}\) defined by \(\lambda \mapsto \mathrm{R}(u,\lambda)\) is not bounded, so there exists \(\lambda\) in \(\mathbf{C}\) such that \(\| \mathrm{R}(u,\lambda)\| >\varepsilon^{-1}\) .

PROPOSITION 19. --- Let E be a complex Banach space and u a closed partial operator on E. Then we have \(\mathrm{PSp}_{\varepsilon}(u) \subset \mathrm{PSp}_{\delta}(u)\) if \(0 < \varepsilon < \delta\) and

\[
\bigcap_ {\varepsilon \in \mathbf {R} _ {+} ^ {*}} \mathrm{PSp} _ {\varepsilon} (u) = \mathrm{Sp} (u).
\]

The first assertion follows from the definition. For the second, the spectrum of u is contained in \(\mathrm{PSp}_{\varepsilon}(u)\) for all \(\varepsilon > 0\) by definition, and if \(\lambda \notin \mathrm{Sp}(u)\) , then \(\lambda \notin \mathrm{PSp}_{\varepsilon}(u)\) when \(\varepsilon < \|R(u, \lambda)\|^{-1}\) .

PROPOSITION 20. --- Let E be a complex Banach space and u a closed partial operator on E. Let \(\varepsilon\in\mathbf{R}_{+}^{*}\) . For every \(\lambda\in\mathbf{C}\) the following conditions are equivalent:

\begin{enumerate}
\def\labelenumi{(\roman{enumi})}
\item
  One has \(\lambda \in \mathrm{PSp}_{\varepsilon}(u)\) ;
\item
  Either \(\lambda \in \mathrm{Sp}(u)\), or there exists \(x \in \mathrm{dom}(u)\) such that \(\| x \| = 1\) and \(\| (\lambda 1_{\mathrm{E}} - u)(x) \| < \varepsilon\) ;
\item
  There exists \(v \in \mathcal{L}(\mathrm{E})\) such that \(\| v \| < \varepsilon\) and \(\lambda \in \mathrm{Sp}(u + v)\) .
\end{enumerate}

We may assume that E is nonzero. Condition (ii) is a reformulation of the definition, and hence of condition (i).

Suppose condition (i) holds, and let us prove (iii). If \(\lambda\) belongs to the spectrum of \(u\) one may take \(v = 0\) in (iii).

Suppose therefore that \(\lambda \notin \mathrm{Sp}(u)\) . By definition of \(\mathrm{PSp}_{\varepsilon}(u)\) , there exists \(y \in \mathrm{E}\) such that \(\| y \| = 1\) and \(\| \mathrm{R}(u, \lambda)y \| > \varepsilon^{-1}\) . Set \(x = \mathrm{R}(u, \lambda)y\) . We have \(x \neq 0\) . By the Hahn-Banach theorem (EVT, II, p. 67, cor. 2), the linear form \(\ell\) on \(\mathbf{C}x\) such that \(\ell(x) = 1\) admits a continuous extension \(\ell_1 \in \mathrm{E}'\) such that \(\| \ell_1 \| = \| \ell \|\) ; we therefore have \(\| \ell_1 \| = \| x \|^{-1} < \varepsilon\) . For every \(e \in \mathrm{E}\) set \(v(e) = \ell_1(e)y\) . We have \(v \in \mathcal{L}(\mathrm{E})\) and \(v(x) = y\) . We obtain \((u + v)x = u(x) + y = \lambda x\) , hence \(\lambda \in \mathrm{Sp}(u + v)\) . Since moreover \(\| v \| = \| \ell_1 \| < \varepsilon\) the condition (iii) is satisfied.

Conversely, let \(v \in \mathcal{L}(\mathrm{E})\) such that \(\|v\| < \varepsilon\) and \(\lambda \in \mathrm{Sp}(u + v)\) . The partial operator \(\lambda1_{\mathrm{E}} - (u + v)\) is therefore either not injective or does not have a continuous inverse; by Lemma 4 of IV, p. 243, applied to the injective partial operator \(\lambda1_{E} - u\) and to \(-v\), we therefore have \(\|v\| \geqslant \|R(u, \lambda)\|^{-1}\) . It follows that condition (iii) implies (i).

COROLLARY. — Let E be a complex Banach space and u a closed partial operator on E. Let \(\varepsilon > 0\) .

\begin{enumerate}
\def\labelenumi{\alph{enumi})}
\item
  For every \(v \in \mathcal{L}(\mathrm{E})\) , we have \(\mathrm{PSp}_{\varepsilon}(u) \subset \mathrm{PSp}_{\varepsilon + \| v\|}(u + v)\) ;
\item
  Let \(\delta > 0\) and let \(D_{\delta}\) be the open disk with center 0 and radius \(\delta\) in \(\mathbf{C}\) . We have \(\mathrm{PSp}_{\varepsilon}(u) + D_{\delta} \subset \mathrm{PSp}_{\varepsilon + \delta}(u)\) .
\end{enumerate}

Let \(\lambda \in \mathrm{PSp}_{\varepsilon}(u)\) . There exists an endomorphism \(w \in \mathcal{L}(\mathrm{E})\) such that \(\|w\| < \varepsilon\) and \(\lambda \in \mathrm{Sp}(u+w)\) (prop. 20). Since \(u+w = (u+v)+(w-v)\) and \(\|w-v\| < \varepsilon + \|v\|\) , we have \(\lambda \in \mathrm{PSp}_{\varepsilon+ \|v\|}(u+v)\) (loc. cit.).

Let \(\mu \in \mathrm{D}_{\delta}\) ; we have \(\lambda + \mu \in \mathrm{Sp}(u + (w + \mu 1_{\mathrm{E}}))\) and \(\| w + \mu 1_{\mathrm{E}} \| < \varepsilon + \delta\) , hence \(\lambda + \mu \in \mathrm{PSp}_{\varepsilon + \delta}(u)\) (loc. cit.).

PROPOSITION 21. --- Let E be a complex Banach space and u a closed partial operator on E. Let \(\varepsilon\in\mathbf{R}_{+}^{*}\) . Every bounded connected component of \(\mathrm{PSp}_{\varepsilon}(u)\) meets the spectrum of \(u\) .

Let U be a connected component of \(\mathrm{PSp}_{\varepsilon}(u)\) which does not meet \(\mathrm{Sp}(u)\) . The set U is open and closed in \(\mathrm{PSp}_{\varepsilon}(u)\) and its closure \(\overline{U}\) in C therefore satisfies \(\overline{U} \cap \mathrm{PSp}_{\varepsilon}(u) = U\) .

Since \(\mathrm{Sp}(u)\) is contained in \(\mathrm{PSp}_{\varepsilon}(u)\) and U does not meet \(\mathrm{Sp}(u)\) , this equality proves that the set \(\overline{U}\) is disjoint from \(\mathrm{Sp}(u)\) , hence contained in the resolvent set of u. Moreover, it implies that the set \(\overline{U}-U\) does not meet \(\mathrm{PSp}_{\varepsilon}(u)\) .

We therefore have \(\|\mathrm{R}(u,\lambda)\| \leqslant \varepsilon^{-1}\) for all \(\lambda\) in \(\overline{U} - U\) , while \(\|\mathrm{R}(u,\lambda)\| > \varepsilon^{-1}\) for \(\lambda \in U\) . If the set U is bounded, its closure \(\overline{U}\) is compact and there exists \(\lambda_{0} \in \overline{U}\) such that \(\|\mathrm{R}(u,\lambda)\| \leqslant \|\mathrm{R}(u,\lambda_{0})\|\) for \(\lambda \in \overline{U}\) . The foregoing implies that \(\lambda_{0} \in U\) which contradicts the maximum principle (VAR, R1, p. 29, 3.3.7) since the resolvent of u is holomorphic in the resolvent set of u (prop. 14 of IV, p. 246).

